% =========================================================
% Completeness and Its Consequences
% =========================================================

\subsection{Completeness of the Real Numbers}

% ---------------------------------------------------------
% Toolkit
% ---------------------------------------------------------
\begin{tcolorbox}[colback=gray!6, colframe=gray!40, arc=2pt,
  left=6pt, right=6pt, top=4pt, bottom=4pt,
  title={\small\textbf{Completeness — Quick Reference}},
  fonttitle=\small\bfseries]
\renewcommand{\arraystretch}{1.4}
\begin{tabular}{@{}p{0.30\textwidth} p{0.44\textwidth} p{0.18\textwidth}@{}}
\textbf{Result} & \textbf{Statement} & \textbf{Proof} \\
\midrule
Completeness Axiom
  & Every nonempty bounded-above $S \subseteq \mathbb{R}$ has $\sup S \in \mathbb{R}$
  & axiom \\
Nested Interval Property
  & $\bigcap_{n=1}^\infty [a_n,b_n] \neq \varnothing$ for nested closed intervals
  & \hyperref[prf:nested-interval-property]{Proof $\to$} \\
Archimedean Property
  & $\forall x \in \mathbb{R}\; \exists n \in \mathbb{N}\; (n > x)$
  & \hyperref[prf:archimedean-property]{Proof $\to$} \\
Floor Lemma
  & $\forall x\; \exists!\, m \in \mathbb{Z}\; (m \le x < m+1)$
  & \hyperref[prf:floor-lemma]{Proof $\to$} \\
Density of $\mathbb{Q}$
  & Between any two reals lies a rational
  & \hyperref[prf:density-of-rationals]{Proof $\to$} \\
Existence of $\sqrt{a}$
  & $\forall a \ge 0\; \exists!\, x \ge 0\; (x^2 = a)$
  & \hyperref[prf:existence-of-square-roots]{Proof $\to$} \\
\end{tabular}
\end{tcolorbox}


% ---------------------------------------------------------
\subsubsection{Completeness Axiom}
% ---------------------------------------------------------

\begin{remark*}[Why completeness is needed]
The ordered field axioms alone do not prevent ``holes'' in the number line.
The rationals $\mathbb{Q}$ satisfy all field and order axioms, yet fail
completeness. Consider the set
\[
S = \{x \in \mathbb{Q} : x^2 < 2\}.
\]
This set is nonempty (e.g.\ $1 \in S$) and bounded above in $\mathbb{Q}$
(e.g.\ $2$ is an upper bound). Yet $\sup S$ does not exist as a rational
number: the candidate $\sqrt{2}$ is irrational. The set $S$ has no least
upper bound in $\mathbb{Q}$ --- a hole sits exactly where $\sqrt{2}$ should be.

The Completeness Axiom asserts that $\mathbb{R}$ has no such holes:
every nonempty bounded-above set has a supremum \emph{inside} $\mathbb{R}$.
This single axiom is what distinguishes $\mathbb{R}$ from $\mathbb{Q}$,
and it underlies every major theorem in analysis.
\end{remark*}

\begin{tcolorbox}[colback=axiombox, colframe=axiomborder, arc=2pt,
  left=6pt, right=6pt, top=4pt, bottom=4pt,
  title={\small\textbf{Axiom (Completeness Axiom (Least Upper Bound Property))}},
  fonttitle=\small\bfseries]
Every nonempty subset of $\mathbb{R}$ that is bounded above has a supremum
in $\mathbb{R}$.

Equivalently: if $S \subseteq \mathbb{R}$ is nonempty and bounded above,
then $\sup S$ exists as a real number.
\end{tcolorbox}

\begin{remark*}[Logical form]
\[
\forall S\;\Bigl(
(S \subseteq \mathbb{R} \wedge S \neq \varnothing \wedge
\exists M \in \mathbb{R}\ \forall x \in S\ (x \le M))
\rightarrow
\exists s \in \mathbb{R}\ (s = \sup S)
\Bigr).
\]
\end{remark*}

\begin{remark*}[Equivalent formulations]
Over the ordered-field axioms, completeness is equivalent to each of the
following:
\begin{itemize}
  \item Every nonempty set bounded below has an infimum.
  \item Every Cauchy sequence in $\mathbb{R}$ converges.
  \item \emph{(Nested Interval Property)} Every nested sequence of nonempty
        closed bounded intervals has nonempty intersection.
\end{itemize}
\end{remark*}

% =========================================================
% Nested Interval Property
% =========================================================

\subsubsection{Nested Interval Property}

\begin{theorem}[\texorpdfstring{\hyperref[prf:nested-interval-property]{Nested Interval Property}}{Nested Interval Property}]
\label{thm:nested-interval-property}
Let $\{[a_n, b_n]\}_{n \in \mathbb{N}}$ be nonempty closed intervals such that
\[
[a_{n+1}, b_{n+1}] \subseteq [a_n, b_n]
\quad \text{for all } n.
\]
Then
\[
\bigcap_{n=1}^\infty [a_n, b_n] \neq \varnothing.
\]
If additionally $b_n - a_n \to 0$, the intersection consists of exactly one
point.
\end{theorem}

\begin{figure}[h]
\centering
\begin{tikzpicture}[x=0.95cm, y=0.9cm, >=Latex]
  % Draw four nested intervals, each row shifted down
  \foreach \row/\la/\rb/\label in {
    0/0.0/10.0/{$[a_1, b_1]$},
    1/1.2/8.5/{$[a_2, b_2]$},
    2/2.0/7.0/{$[a_3, b_3]$},
    3/2.6/5.8/{$[a_4, b_4]$}
  }{
    \draw[thick, blue!70!black] (\la, -\row*0.8) -- (\rb, -\row*0.8);
    \fill[blue!70!black] (\la, -\row*0.8) circle (2.0pt);
    \fill[blue!70!black] (\rb, -\row*0.8) circle (2.0pt);
    \node[left] at (\la-0.1, -\row*0.8) {\small \label};
  }
  % Indicate left endpoints climbing
  \draw[->, dashed, red!70!black, thick] (0.0,-0.0) -- (0.0,-3.2);
  \node[right, red!70!black] at (0.05, -1.6) {\small $(a_n)\!\uparrow$};
  % Indicate right endpoints falling
  \draw[->, dashed, red!70!black, thick] (10.0,0.0) -- (10.0,-3.2);
  \node[left, red!70!black] at (9.95, -1.6) {\small $(b_n)\!\downarrow$};
  % Pinch point
  \draw[dotted] (4.0, 0.4) -- (4.0, -3.3);
  \node[above] at (4.0, 0.45) {\small $x \in \bigcap [a_n,b_n]$};
\end{tikzpicture}
\caption{Nested intervals $[a_1,b_1] \supseteq [a_2,b_2] \supseteq \cdots$ pinch toward a
common point. The sequence of left endpoints is increasing and bounded above; its
supremum $x$ lies in every interval.}
\label{fig:nested-intervals}
\end{figure}



% =========================================================
% Archimedean Property
% =========================================================

\subsubsection{Archimedean Property}

\begin{remark*}[Why this needs a proof]
The Archimedean property asserts that the natural numbers are unbounded in $\mathbb{R}$.
This might seem obvious --- after all, $1, 2, 3, \ldots$ keep growing --- but
that intuition assumes we have already constructed $\mathbb{N}$ inside $\mathbb{R}$
from first principles. Here we are working axiomatically, and the ordered-field
axioms alone do not prevent $\mathbb{N}$ from being bounded. (Non-Archimedean
ordered fields, such as fields of formal Laurent series, satisfy all ordered-field
axioms yet have bounded $\mathbb{N}$.)

The Completeness Axiom rules this out: if $\mathbb{N}$ were bounded above in
$\mathbb{R}$, it would have a supremum $\alpha \in \mathbb{R}$ by the LUB
property. But then $\alpha - 1$ is not an upper bound (by the
$\varepsilon$-characterization with $\varepsilon = 1$), so some $n \in \mathbb{N}$
satisfies $n > \alpha - 1$, giving $n + 1 > \alpha$ --- a natural number larger
than the supposed supremum.
\end{remark*}

\begin{tcolorbox}[colback=propbox, colframe=propborder, arc=2pt,
  left=6pt, right=6pt, top=4pt, bottom=4pt,
  title={\small\textbf{Definition}},
  fonttitle=\small\bfseries]
$\mathbb{R}$ satisfies the \emph{Archimedean property} if
\[
\forall x \in \mathbb{R}\; \exists n \in \mathbb{N}\; (n > x).
\]
Equivalently:
\[
\forall x > 0\; \forall y \in \mathbb{R}\; \exists n \in \mathbb{N}\; (nx > y).
\]
\end{tcolorbox}

\begin{remark*}[Logical form]
\[
\forall x\; \exists n\; (n > x).
\qquad
\forall x > 0\; \forall y\; \exists n\; (nx > y).
\]
\end{remark*}

\begin{theorem}[\texorpdfstring{\hyperref[prf:archimedean-property]{Archimedean Property}}{Archimedean Property}]
\label{thm:archimedean-property}
$\mathbb{R}$ satisfies the Archimedean property.
\end{theorem}



\begin{corollary}[Archimedean Property — Corollary]
If $x > 0$ and $y \in \mathbb{R}$, then $\exists n \in \mathbb{N}$ such that
$nx > y$.
\end{corollary}



% =========================================================
% Integer Part
% =========================================================

\subsubsection{Integer Part}

\begin{lemma}[\texorpdfstring{\hyperref[prf:floor-lemma]{Floor Lemma}}{Floor Lemma}]
\label{lem:floor-lemma}
For every $x \in \mathbb{R}$ there exists a unique $m \in \mathbb{Z}$ such that
\[
m \le x < m + 1.
\]
\end{lemma}

\begin{remark*}[Worked example]
Before proving the lemma, verify it on two concrete cases.

\medskip
\textbf{Case $x = 2.7$:} we seek $m \in \mathbb{Z}$ with $m \le 2.7 < m+1$.
Taking $m = 2$ gives $2 \le 2.7 < 3$. \checkmark

\medskip
\textbf{Case $x = -1.3$:} we seek $m \le -1.3 < m+1$.
One might guess $m = -1$, but $-1.3 < -1$, so $-1 \le -1.3$ fails.
The correct choice is $m = -2$: $-2 \le -1.3 < -1$. \checkmark

The key point for negative non-integers is that the floor rounds \emph{down}
(toward $-\infty$), not toward zero.  The proof below constructs $m$ via the
Archimedean property, which applies uniformly to positive and negative $x$.
\end{remark*}



% =========================================================
% Density
% =========================================================

\subsubsection{Density}

\begin{tcolorbox}[colback=propbox, colframe=propborder, arc=2pt,
  left=6pt, right=6pt, top=4pt, bottom=4pt,
  title={\small\textbf{Definition}},
  fonttitle=\small\bfseries]
A subset $A$ of a linearly ordered set $X$ is \emph{dense in $X$} if
\[
\forall a < b\; \exists c \in A\; (a < c < b).
\]
\end{tcolorbox}

\begin{remark*}[Logical form for $\mathbb{Q}$ dense in $\mathbb{R}$]
\[
\forall a < b\; \exists q \in \mathbb{Q}\; (a < q < b).
\]
\end{remark*}

\begin{theorem}[\texorpdfstring{\hyperref[prf:density-of-rationals]{Density of $\mathbb{Q}$}}{Density of Q}]
\label{thm:density-of-rationals}
$\mathbb{Q}$ is dense in $\mathbb{R}$.
\end{theorem}



\begin{corollary}[Irrationals Are Dense]
Between any two distinct real numbers lies an irrational.
\end{corollary}

\begin{remark*}[Proof idea: scale to hunt for a rational, then multiply back]
We want an irrational $t \in (a, b)$.
Irrationals are harder to construct directly, so the strategy is:

\medskip
\textbf{Observation:} if $q$ is a nonzero rational and $\sqrt{2}$ is irrational,
then $q\sqrt{2}$ is irrational.

\medskip
\textbf{Plan:} instead of hunting for an irrational in $(a, b)$, hunt for a
rational $q \in (a/\sqrt{2},\, b/\sqrt{2})$ (possible by Density of $\mathbb{Q}$),
then set $t = q\sqrt{2}$.
Since $a/\sqrt{2} < q < b/\sqrt{2}$, multiplying through by $\sqrt{2} > 0$ gives
$a < t < b$, and $t = q\sqrt{2}$ is irrational.
\end{remark*}

\begin{corollary}[Density of Irrationals]
\label{cor:density-of-irrationals}
The irrationals are dense in $\mathbb{R}$.
\end{corollary}

% =========================================================
% Existence of Square Roots
% =========================================================

\subsubsection{Existence of Square Roots}

\begin{theorem}[\texorpdfstring{\hyperref[prf:existence-of-square-roots]{Existence of Square Roots}}{Existence of Square Roots}]
\label{thm:existence-of-square-roots}
For every $a \ge 0$ there exists a unique $x \ge 0$ such that $x^2 = a$.
\end{theorem}

\begin{remark*}[Logical form]
\[
\forall a \ge 0\; \exists!\, x \ge 0\; (x^2 = a).
\]
\end{remark*}





% =========================================================
% Structural Summary
% =========================================================

\subsubsection*{Structural Summary}

\[
\text{Field Axioms}
\;\Rightarrow\;
\text{Order Axioms}
\;\Rightarrow\;
\text{Completeness Axiom}
\]

Completeness yields, in logical order:

\[
\text{Nested Interval Property}
\;\Rightarrow\;
\text{Archimedean Property}
\;\Rightarrow\;
\text{Floor Lemma}
\;\Rightarrow\;
\text{Density of } \mathbb{Q}
\;\Rightarrow\;
\text{Existence of } \sqrt{a}
\]

Completeness is the property that prevents holes in $\mathbb{R}$. It is
equivalent to the Cauchy Criterion and underlies every major limit theorem
in real analysis.

\begin{figure}[h]
\centering
\begin{tikzpicture}[
  node distance=0.5cm and 0.6cm,
  box/.style={draw, rounded corners=3pt, fill=gray!8, minimum width=3.1cm,
               minimum height=0.65cm, align=center, font=\small},
  axiombox/.style={draw, rounded corners=3pt, fill=axiombox,
               draw=axiomborder, minimum width=3.1cm,
               minimum height=0.65cm, align=center, font=\small},
  arr/.style={->, thick, >=Latex}
]
  \node[axiombox] (comp) {Completeness Axiom};

  \node[box, below left=0.9cm and 1.5cm of comp]  (nip)   {Nested Interval\\Property};
  \node[box, below right=0.9cm and 1.5cm of comp] (sqrt)  {Existence of $\sqrt{a}$};

  \node[box, below=0.9cm of nip]   (arch)  {Archimedean Property};
  \node[box, below=0.9cm of arch]  (floor) {Floor Lemma};
  \node[box, below=0.9cm of floor] (densQ) {Density of $\mathbb{Q}$};
  \node[box, below=0.9cm of densQ] (densI) {Density of Irrationals};

  \draw[arr] (comp) -- (nip);
  \draw[arr] (comp) -- (sqrt);
  \draw[arr] (nip)  -- (arch);
  \draw[arr] (arch) -- (floor);
  \draw[arr] (floor)-- (densQ);
  \draw[arr] (densQ)-- (densI);
\end{tikzpicture}
\caption{Logical dependency chain from the Completeness Axiom.
Every result in this chapter ultimately flows from the single axiom at the top.}
\label{fig:completeness-chain}
\end{figure}

\begin{remark*}[Connections forward]
Completeness appears in every major limit theorem in the notes.
The Monotone Convergence Theorem (convergence chapter) constructs the
limit of a bounded monotone sequence as $\sup\{a_n : n \in \mathbb{N}\}$;
existence of this supremum requires completeness.
The Bolzano--Weierstrass Theorem applies the Nested Interval Property
directly to extract a convergent subsequence from any bounded sequence.
The Cauchy Criterion identifies convergent sequences purely in terms of
their own spacing, without naming the limit in advance; its proof
constructs the limit as a supremum and again requires completeness.
In metric spaces (Distance chapter), completeness in the sense of
Cauchy sequences is the natural generalisation of the same idea to
abstract settings.
Every $\varepsilon$-$\delta$ proof that locates a limit or extremum is
drawing on this chapter's infrastructure.
\end{remark*}